% ---------------------------------------------------------------------------------------
% MODELO DE TCC - INSTITUTO FEDERAL DE SERGIPE (IFS)
% ---------------------------------------------------------------------------------------
% Autor do Modelo: Mikael dos Santos Gois (mikaelgois1@gmail.com)
% Baseado nos modelos de: 
%   - Diogo D. Moreira (TCC, Dissertação e Tese)
%   - DECOM/UFOP (Monografia de Ciência da Computação)
% Modelo Base: abnTeX2 adaptado para as normativas internas de TCC do IFS (Art. 31 e 34)
% Versão: 1.0 (2026)
% ---------------------------------------------------------------------------------------
% Este arquivo faz parte de um modelo opensource para estudantes do IFS.
% Não remova os créditos dos autores originais e do modelo base.
% Este modelo é fornecido "no estado em que se encontra", sem garantias de qualquer tipo.
% ---------------------------------------------------------------------------------------

% Anexo:
% Elemento opcional. O texto deve ser precedido da palavra ANEXO, identificado por letras maiúsculas
% consecutivas, travessão e pelo respectivo título. Na identificação dos anexos, quando esgotadas as
% letras do alfabeto, devem ser utilizadas letras maiúsculas dobradas. O destaque tipográfico (negrito,
% itálico ou sublinhado) deve ser o mesmo da seção primária. - NBR 14724:2024 (4.2.3.4)

% EXEMPLO
% ANEXO A — REPRESENTAÇÃO GRÁFICA DE CONTAGEM DE CÉLULAS INFLAMATÓRIAS PRESENTES NAS CAUDAS EM REGENERAÇÃO – GRUPO DE CONTROLE I (TEMPERATURA...)

\begin{anexosenv}
  % \partanexos   % indica o início dos anexos

  % [O que aparece no PDF Bookmark] {O que aparece impresso no documento}
  \chapter[Exemplo]{EXEMPLO} \label{anex:a}
  Utiliza a mesma formatação dos apêndices.
\end{anexosenv}